% ===================================================================
%  TABELL Nr 11 — Staden och dess minnesmärken. — Branden.
%  Traduction 2026 d'apres texto/io, controlee sur texto/fr et
%  texto/es. Consigne : voir texto/sv/10-tabell-01.tex.
% ===================================================================

%%K t11-apar-1 apar
\VUsaut{2.0mm}
\VUcentre{13.2pt}{90}{TREDJE SERIEN}
\VUsaut{3.0mm}
\VUfilet{16mm}
\VUsaut{5.6mm}
\VUtitre{15.0pt}{60}{TABELL Nr 11}
\VUsaut{1.4mm}
\VUpk{11.4pt}{40}{Staden och dess minnesmärken. --- Branden.}
\VUsaut{2.2mm}

%%K t11-tit-1 sub
\VUpk{10.2pt}{40}{Förstaden.}
\VUsaut{3.0mm}

%%K t11-c1-01-1 p
1. --- Jag bor i en stor stad som ligger 100 kilometer från Oceanen, och som ändå är en
\VUgras{hamn} \textsuperscript{(1)} av första rangen. Floden som gränsar till den har en bredd
av 400 meter och bildar en mycket vacker redd.

%%K t11-c1-02-1 p
2. --- Hela den del av staden som ligger vid \VUgras{kajerna} \textsuperscript{(2)} ser helt
sjöfarts- och industribetonad ut, och bildar en
\VUgras{förstad} \textsuperscript{(3)} som bebos bara av sjöfolk och arbetare.
Dessa arbetar i \VUgras{fabrikerna} \textsuperscript{(4)} och i
\VUgras{gasverket} \textsuperscript{(5)} vars höga \VUgras{skorsten}
\textsuperscript{(6)} och \VUgras{gasklocka} syns på mycket långt håll.

%%K t11-c1-03-1 p
3. --- Under medeltiden var staden befäst, och man finner ännu några rester av dess vallar, i
synnerhet \VUgras{porten} \textsuperscript{(8)} till det gamla
\VUgras{befästa slottet} \textsuperscript{(7)} flankerad av sina två
\VUgras{torn} \textsuperscript{(9)} som nu tjänar som \VUgras{fängelse}.

%%K t11-c1-04-1 p
4. --- I hela den \VUgras{stadsdelen}, som ser mycket gammal ut, är bara en byggnad
förhållandevis modern: det är \VUgras{tullhuset} \textsuperscript{(10)}. Det ligger mellan
kajen och den \VUgras{lilla gatan} \textsuperscript{(11)} som leder till det gamla
\VUgras{rådhuset} \textsuperscript{(12)}. Man känner lätt igen det sistnämnda minnesmärket på
det jättelika \VUgras{klocktornet} \textsuperscript{(13)} som behärskar den gamla staden.

%%K t11-c1-05-1 p
5. --- På andra sidan gatan \VUgras{står} en byggnad uppförd i samma stil som tullhuset: det
är \VUgras{Börsen} \textsuperscript{(14)}. En av dess fasader vetter mot ett litet torg där
\VUgras{länsstyrelsen} \textsuperscript{(15)} ligger. Vår stad är nämligen, utan att vara en
huvudstad, ändå en viktig länshuvudstad.

%%K t11-tit-2 sub
\VUsaut{5.0mm}
\VUpk{10.2pt}{40}{Stadens högsta del.}
\VUsaut{3.4mm}

%%K t11-c1-06-1 p
6. --- Garnisonen, ganska talrik, är inkvarterad i vidsträckta mycket sunda
\VUgras{kaserner} \textsuperscript{(16)}; de sträcker sig ända till
\VUgras{stadsparkens} \textsuperscript{(17)} galler. \VUgras{Boulevarden}
\textsuperscript{(19)} som leder till den parken går bakom \VUgras{sparbanken}
\textsuperscript{(18)} och slutar vid \VUgras{katedralens} \textsuperscript{(20)} torg.

%%K t11-c1-07-1 p
7. --- Alla de minnesmärkena sträcker sig som en amfiteater på en kulle där också vårt
vidsträckta \VUgras{läroverk} \textsuperscript{(21)} ligger.

%%K t11-c1-07-2 p
Om man går ner längs en något sluttande väg, som möter Storgatan, kommer man till
\VUgras{Domstolen} \textsuperscript{(22)} framför vilken en behaglig
\VUgras{offentlig trädgård} \textsuperscript{(26)} sträcker sig. Den
\VUgras{planteringen} är inte särskilt stor, men den bildar mitt i staden en oas
av grönska och svalka. Därför besöks den mycket av stadsborna som gärna kommer och sätter sig
under \VUgras{kaféets} \textsuperscript{(25)} tälttak, för att lyssna på soldatmusikanterna
som på torsdagar och söndagar spelar i
\VUgras{kiosken} \textsuperscript{(27)} som står mellan gräsmattorna och
buskagen.

%%K t11-c1-08-1 p
8. --- På en av de gräsmattorna står en \VUgras{staty} \textsuperscript{(28)} av brons, ett
minnesmärke rest till minne av en av våra stadsbor, som gjorde stora tjänster åt sin
födelsestad.

%%K t11-tit-3 sub
\VUsaut{4.0mm}
\VUpk{10.2pt}{40}{Transportmedlen.}
\VUsaut{3.4mm}

%%K t11-c1-09-1 p
9. --- Hittills är vår stad ännu inte upplyst med elektriskt ljus, den har bara
\VUgras{gaslyktor} \textsuperscript{(29)}. Men \VUgras{ortstidningarna} för en
kampanj för att det belysningssystemet skulle förbättras, \VUgras{liksom} man redan har
fulländat \VUgras{spårvagnarnas dragningssystem}. \VUgras{De} gamla vagnarna, \VUgras{dragna}
av \VUgras{hästar}, har praktiskt taget ersatts av
\VUgras{elektriska spårvagnar} \textsuperscript{(31)} som snabbt för de resande
från \VUgras{stadens} ena ände \VUgras{till den andra}, till ett mycket
\VUgras{lågt} pris.

%%K t11-c1-10-1 p
10. --- Vid Domstolen ligger \VUgras{Banken} \textsuperscript{(32)} där man diskonterar
handelspappren. \VUgras{Banktjänstemännen} \textsuperscript{(33)} går och inkasserar skulderna
i hemmen.

%%K t11-tit-4 sub
\VUsaut{4.0mm}
\VUpk{10.2pt}{40}{Konst och undervisning.}
\VUsaut{3.4mm}

%%K t11-c1-11-1 p
11. --- Den vackra byggnaden som står nära intill, är \VUgras{museet}. Det innehåller på samma
gång stadens \VUgras{bibliotek} \textsuperscript{(34)}, de vackra
\VUgras{naturvetenskapliga samlingarna} \textsuperscript{(35)} (Naturhistoriska museet) och
ett prydligt \VUgras{tavelmuseum} \textsuperscript{(36)} som utgör en praktfull ständig
\VUgras{utställning}.

%%K t11-c1-11-2 p
Det öppnas för allmänheten två gånger i veckan, men främlingarna får besöka det vilken dag som
helst mot en drickspenning åt \VUgras{vaktmästaren} \textsuperscript{(37)}. \VUgras{Målarna}
\textsuperscript{(38)} kommer dit för att arbeta. Man ser dem \VUgras{framför} sitt staffli,
med en \VUgras{palett} i handen, kopiera de berömda tavlorna.

%%K t11-c1-12-1 p
12. --- På höjden, bakom museet, börjar vi se \VUgras{kupolen} \textsuperscript{(40)} på
\VUgras{sjukhuset} \textsuperscript{(39)} och byggnaderna till \VUgras{seminariet}
\textsuperscript{(41)} i vilket lärarna vid våra kommunala skolor utbildades. På Teatertorget
står ett \VUgras{postkontor} \textsuperscript{(43)} inrymt i ett \VUgras{privat hus}
\textsuperscript{(42)}.

%%K t11-tit-5 sub
\VUsaut{5.0mm}
\VUpk{11.4pt}{40}{Branden.}
\VUsaut{3.0mm}

%%K t11-tit-6 sub
\VUpk{10.2pt}{40}{Brandrop.}
\VUsaut{3.4mm}

%%K t11-c2-01-1 p
1. --- Ett förskräckligt rykte sprider sig genom staden: det har just brutit ut en brand på
teatern! I det ögonblick då jag kommer fram till \VUgras{torget} \textsuperscript{(96)}, står
folkmassan redan samlad på \VUgras{trottoaren} \textsuperscript{(46)} och på \VUgras{körbanan}
\textsuperscript{(47)}.
\VUgras{Posttjänstemännen} \textsuperscript{(44)} och \VUgras{brevbärarna}
\textsuperscript{(45)} kommer ut ur postkontoret. En \VUgras{spårvagnsförare}
\textsuperscript{(48)} måste stanna sin (spår)vagn för att låta en kommunal
\VUgras{ambulansvagn} \textsuperscript{(50)} passera, som kommer med en
\VUgras{kirurg} \textsuperscript{(52)} och en \VUgras{sjukvårdare}
\textsuperscript{(51)} för att sköta en \VUgras{sårad} \textsuperscript{(53)} som burits på en
\VUgras{bår} \textsuperscript{(54)} till
\VUgras{tidningskiosken} \textsuperscript{(55)}.

%%K t11-c2-02-1 p
2. --- Ett infanterikompani, som kom tillbaka från övningarna, stannade för att säkra
ordningstjänsten tillsammans med den \VUgras{beridna stadspolisen} \textsuperscript{(61)}.
\VUgras{Ryttarna}, vilkas hästar stegrar sig, får, bättre än \VUgras{soldaterna}
\textsuperscript{(57)} eller \VUgras{gendarmerna} \textsuperscript{(63)}, \VUgras{de nyfikna}
\textsuperscript{(56)} att vika tillbaka. För övrigt har gendarmerna mycket att göra. En av
dem visar
\VUgras{poliskonstaplarna} \textsuperscript{(64)} vart man ska föra en annan
\VUgras{skadad} \textsuperscript{(65)}, en modig brandman som har fallit ner från
en stege.

%%K t11-c2-03-1 p
3. --- \VUgras{Officeren} \textsuperscript{(66)} anger noga den plats där man ska ställa
\VUgras{ångsprutan} \textsuperscript{(67)} som är på väg. I väntan på den kraftiga maskinen
har han satt upp en \VUgras{handspruta} \textsuperscript{(68)} vars långa \VUgras{slang}
\textsuperscript{(69)} kastar vattnet som en störtflod på elden. Sex brandmän manövrerar den,
medan deras kamrat, som håller \VUgras{strålröret} \textsuperscript{(70)}, riktar
\VUgras{strålen} \textsuperscript{(71)} mot brandens medelpunkt.

%%K t11-c2-04-1 p
4. --- På andra sidan teatern har man rest den stora \VUgras{stegen} \textsuperscript{(72)}.
Den gör det möjligt för brandmännen att närma sig lågorna som slår upp mellan virvlar av svart
\VUgras{rök} \textsuperscript{(75)}.

%%K t11-tit-7 sub
\VUsaut{5.0mm}
\VUpk{10.2pt}{40}{Tappra bragder.}
\VUsaut{3.4mm}

%%K t11-c2-05-1 p
5. --- \VUgras{Branden} \textsuperscript{(74)} uppstod i \VUgras{teaterns}
\textsuperscript{(73)} foajé, medan man repeterade \VUgras{operan} vars första
\VUgras{föreställning} skulle ha ägt rum i morgon. Lyckligtvis var det bara få
\VUgras{åskådare} \textsuperscript{(76)} i salongen. De olyckliga har flytt till
\VUgras{balkongen} \textsuperscript{(77)} där modiga \VUgras{brandmän}
\textsuperscript{(86)} kommer för att hjälpa dem med stegar och tåg.

%%K t11-c2-06-1 p
6. --- Under \VUgras{peristylen} \textsuperscript{(78)}, där \VUgras{affischen}
\textsuperscript{(79)} upplyser om föreställningen, ser man \VUgras{musikerna}
\textsuperscript{(81)} i \VUgras{orkestern} \textsuperscript{(80)} fly, medan de bär sina
instrument: \VUgras{fiol} \textsuperscript{(81)}, \VUgras{flöjt} \textsuperscript{(82)},
\VUgras{violoncell} \textsuperscript{(83)},
\VUgras{trombon} \textsuperscript{(84)}, \VUgras{trumma}
\textsuperscript{(85)}, o. s. v.. Man söker rädda en del av \VUgras{möblemanget}
\textsuperscript{(87)}.

%%K t11-c2-07-1 p
7. --- \VUgras{Skådespelarna} \textsuperscript{(89)} och
\VUgras{skådespelerskorna} \textsuperscript{(88)}, \VUgras{sångarna} och
\VUgras{sångerskorna} måste fly i sin \VUgras{teaterdräkt}
\textsuperscript{(90)}. Tenoren förklarar för en \VUgras{journalist} \textsuperscript{(92)}
hur det gick till. \VUgras{Reportern} skyndade dit genast från första ögonblicket med
myndighetspersonerna: \VUgras{landshövdingen} \textsuperscript{(91)}, \VUgras{borgmästaren}
\textsuperscript{(93)},
\VUgras{polismästaren} \textsuperscript{(94)} och \VUgras{den rättsliga
tjänstemannen} \textsuperscript{(95)}, som ska undersöka saken för att döma om ansvaret.
\VUsaut{6.0mm}
\VUfilet{22mm}
